\subsection{Cubic incidence and Lorentzian realization}

The cubic completion provides a relation between interior and boundary.
Its geometry can be developed through the incidence of the six faces
with the eight vertices. Label the faces by \((i,\varepsilon)\),
\(i\in\{1,2,3\}\), \(\varepsilon\in\{+1,-1\}\), and the vertices by
\(\nu\in\{+1,-1\}^3\). The incidence matrix is
\[
 M_{(i,\varepsilon),\nu}=\mathbf1_{\{\nu_i=\varepsilon\}}.
\]
Let \(u=(1,\ldots,1)^{\mathsf T}\in\RR^6\), let \(R_F\) denote face
opposition, and let \(P_u=uu^{\mathsf T}/6\), \(P_-=(I-R_F)/2\).

\begin{proposition}[Four-dimensional incidence quotient]
The incidence Gram matrix satisfies
\[
 MM^{\mathsf T}=12P_u+4P_-.
\]
Its eigenvalues \(12,4,0\) have multiplicities \(1,3,2\).
Consequently,
\[
 Q_\square=\RR^6/\ker(MM^{\mathsf T})
 \simeq\RR u\oplus E_-,\qquad E_-=\operatorname{im}P_-.
\]
\end{proposition}
\begin{proof}
A face contains four vertices; two opposite faces have empty intersection,
and two distinct nonopposite faces share two vertices.
The matrix \(12P_u+4P_-\) has precisely those entries.
The uniform space has dimension one; the differences between opposite
faces generate a three-dimensional space orthogonal to the uniform one.
Its complement has dimension two and lies in the kernel.
\end{proof}

The \(1+3\) decomposition arises from incidence. The temporal choice
is specified by the bilinear form
\[
 B_\square([x],[y])=\langle x,(P_--P_u)y\rangle.
\]
It is well defined on the quotient because both projections annihilate the
kernel. In the basis
\[
 t=[u]/\sqrt6,\qquad
 d_i=[e_{i,+}-e_{i,-}]/\sqrt2,
\]
its matrix is \(\operatorname{diag}(-1,1,1,1)\). This makes explicit
the domain, the decomposition, and the sign convention that produce
the Minkowski realization. The positive Gram matrix determines the quotient;
the temporal assignment determines the indefinite form on it.

The operator
\[
 W_\square=\sqrt6P_u+\sqrt2P_-,\qquad
 W_\square e_{i,\varepsilon}=t+\varepsilon d_i
\]
maps the six face labels to null directions, since
\(B_\square(t+\varepsilon d_i,t+\varepsilon d_i)=-1+1=0\).
Direction reversal and face opposition thus acquire a concrete
geometric realization.

\subsubsection{Gluing, reversal, and subdivision}

The cellular realization uses an invertible transport \(U_m(a)\),
which preserves the Lorentzian form, between the fibers at the endpoints of
an oriented dual edge \(a\), a face label
\(\lambda_m(a)\), and an orientation \(\sigma_m(a)\).
With \(\ell_m=\ell_0/9^m\), the gluing map is
\[
 \beta_m(a)=\ell_m\sigma_m(a)
 W_{\square,s(a)}e_{\lambda_m(a)}.
\]
Compatible reversal satisfies
\[
 \beta_m(\bar a)=-U_m(a)\beta_m(a).
\]
Identify the source fibers at two consecutive levels by means of
the specified refinement isometry. If the nine subsegments
\(a_1,\ldots,a_9\), at level \(m+1\), preserve the label
and orientation when transported to that common fiber, then
\begin{equation}
 \beta_m(a)=\sum_{j=1}^9
 U_{m+1}(a_1\cdots a_{j-1})^{-1}\beta_{m+1}(a_j).
 \label{eq:soldadura-refinamiento}
\end{equation}
Each transported term equals \(\ell_m/9\) times the same
null direction, which proves the identity. The compatibility hypothesis
determines which subdivisions realize this refinement.
The equation brings together direction, scale, and transport, instead of reducing
the geometric correspondence to a dimension count.

Once the orientation and Lorentzian form have also been fixed, Hodge duality
on \(\Lambda^2Q_\square^*\) satisfies \(\star^2=-I\).
In dimension four, the sign is \((-1)^{2(4-2)+1}=-1\).
This yields a real complex structure on two-forms and, after
complexification, the eigenspaces with eigenvalues \(+i\) and
\(-i\). The electromagnetic realization will use this domain
after constructing the electronic section and specifying its coupling.

\subsection{Compatibility between modular and positional reductions}

The block-decimal base enters through Euclidean division,
which simultaneously preserves residue and quotient:
\[
 n=1000q+r,\qquad 0\le r<1000.
\]
Since \(1000\equiv1\pmod9\) and \(1000\equiv1\pmod3\),
\[
 n\equiv q+r\pmod9,\qquad n\equiv q+r\pmod3.
\]
Therefore, reduction of a block requires the contribution of the
quotient in order to recover the class of the integer. The numbers \(0\) and
\(1000\) have the same remainder modulo \(1000\), but different classes
modulo \(9\) and modulo \(3\). This is an operational reason to preserve
memory during a change of resolution.

Ternary refinement and determination of a decimal prefix
can be coordinated by means of a single integer residue. If \(N/3^L\)
is a cylinder endpoint and \(D/1000^r\) a decimal endpoint, define
\[
 A=1000^rN-D3^L.
\]
Upon adding six ternary symbols, \(N'=729N+\nu\),
\(L'=L+6\), with \(0\le\nu<729\). If the decimal prefix remains
fixed, then
\[
 A'=729A+\nu1000^r.
\]
If a decimal block \(\delta\) is also incorporated, then
\(D'=1000D+\delta\), \(r'=r+1\), and
\begin{equation}
 A'=729000A+\nu1000^{r+1}-729\delta3^L.
 \label{eq:residuo-dos-bases}
\end{equation}
Both identities are proved by substituting the definitions. Comparison
of endpoints is thus reduced to integer operations that retain the
contribution of each extension. This compatibility is what permits a region
to be refined before its real coordinate is evaluated.
